\documentclass[a4paper,12pt]{article}
\usepackage[utf8]{inputenc}
\usepackage[spanish]{babel}
\usepackage{geometry}
\geometry{margin=1in}
\usepackage{parskip}

\begin{document}

\begin{center}
	\textbf{Carta de Recomendación}
\end{center}

Estimado/a miembro del Comité de Admisión,

Nos complace escribir esta carta de recomendación para Gabriel, quien ha sido nuestro estudiante en la Universidad Politécnica de Madrid en la Escuela Técnica Superior de Ingenieros Informáticos (ETSIINF) durante el Grado en Matemáticas e Informática. Como profesoras de Gabriel en varias asignaturas clave de su carrera, hemos tenido la oportunidad de observar de cerca su desarrollo académico y su dedicación a los estudios matemáticos.

Gabriel cursó con nosotras las asignaturas de Cálculo 1 y Cálculo 2, impartidas por la profesora Raquel, donde demostró un sólido entendimiento de los conceptos fundamentales y una notable capacidad para resolver problemas complejos, obteniendo calificaciones de notable en ambas materias. Posteriormente, en la asignatura de Análisis Funcional, impartida conjuntamente por ambas, Gabriel no solo destacó por su excelente rendimiento académico, obteniendo un sobresaliente, sino también por su interés y curiosidad intelectual, lo que le llevó a profundizar en temas avanzados más allá del currículo estándar.

Fue precisamente durante las clases de Análisis Funcional donde Gabriel se interesó por los polinomios ortogonales de Sobolev, un tema que captó su atención y que, tras varias discusiones con nosotras, decidió explorar en mayor profundidad para su Trabajo de Fin de Grado (TFG). Bajo nuestra supervisión, Gabriel desarrolló un trabajo riguroso y original sobre este tema, demostrando una notable capacidad para la investigación independiente y un profundo entendimiento de los conceptos matemáticos involucrados.

Además, Gabriel ha expresado su intención de continuar esta línea de investigación en su Trabajo de Fin de Máster (TFM), lo cual consideramos una excelente decisión dada su pasión y aptitud para el tema.

Por todas estas razones, recomendamos encarecidamente a Gabriel para su admisión en el programa de máster. Estamos convencidas de que aprovechará al máximo esta oportunidad y continuará destacando en sus estudios y en su futura carrera profesional.

Atentamente,\\
\vspace{1cm}

Raquel Natividad Gonzalo Palomar\\
Profesora del Departamento de Matemática Aplicada a las Tecnologías de la Información y las Comunicaciones\\
Universidad Politécnica de Madrid - ETSIINF\\
\vspace{1cm}

María Del Carmen Escribano Iglesias\\
Profesora del Departamento de Matemática Aplicada a las Tecnologías de la Información y las Comunicaciones\\
Universidad Politécnica de Madrid - ETSIINF\\

\end{document}